\section{Discussion}
\label{sec:discussion}
%% Packet W-10. Ledger rows: NEG.gb5_repair_sweep_not_built. Carries the full nonclaims ledger
%% (design-pack Phase 5 gate: printed in full), the anti-reductionism paragraph, and the gap walk.

\subsection{What is established, at which grade}

Theorem-grade, each on its exactly declared finite class: the energy readout is not closed, with a
certified margin over the equilibrium baseline (GT1); the Kovacs split-pair witness (GT2);
foreclosure of all $512$ energy-definable predicates on the declared class, with a finer
present-state readout as the exhibited separator (GT3); the kill-test
panel with five look-alike regimes correctly rejected (GT4); $\mathrm{MaxFiber}=2$ with a
saturating buyback ledger (GT5); and, from the formalization effort itself, the collapse theorem
for \texttt{RouteTransportCore} (machine-checked). Certificate-grade, shell-local: the GT6 strict
extension via non-factorization. Run-only: the GT7 protocol and its 8/14/294/22 scoreboard.
Recognition-grade: the GB1/GB2 demonstrations, the Lean fidelity filing, and the three empirical
bridges. Nothing in this paper is graded higher than its artifact.

\subsection{What is open}

Seven gaps are disclosed, each with a register entry stating what closing it would take
(\texttt{paper/gap\_register.md}); none is load-bearing for any claim above, because every claim's
wording already excludes it.

\claimref{NEG.gb5\_repair\_sweep\_not\_built}{3f3c56405847a079f48bae350b3b35e901a5731470bcb4a092bd7f682f075559}
The largest is GB5: no exhaustive sweep over a frozen class of candidate repairs exists, so no
repair-impossibility certificate of any grade is claimed. ``No admissible repair exists'' is never
asserted, and GT5's buyback saturation is explicitly not a substitute. The others: the GT2 loop
certificate (unbuilt, and now known to need new formal footing); GT1's [B]-parity control
(skipped, filed); GT1's breadth (one declared (model, $N$, lens) cell, which is the declared-class
discipline working as designed, but honestly a single cell); the GT7 failure interpretation
(three readings left open by design); material forcing (absent on this shell, so [C]'s fuller
route is unexercised); and the numeric-witness Lean encoding (proved impossible in this core).
Natural next experiments follow directly: a second pre-registered prediction round with
re-derived intervals; an Arrhenius-form trap family (per-trap rather than common temperature
dependence) so the Kovacs protocol is non-degenerate there; the FA, $L_{\mathrm{panel}}$, and
$N{=}10$ extensions of GT1; and a weakened formal core admitting numeric witnesses.

One ceiling is declared rather than open, per the framework's recognition-open caveat: if glass
aging carries a finite conservation or balance law over the declared records that supplies none
of the six witness types, the six-role decomposition underlying this paper's filings is
conditional on the recognition hypothesis that no such feature exists. No such candidate was
identified on the audited class; any future one files at recognition grade, separately.

\subsection{Anti-reductionism, stated plainly}

This paper's certificates live at a boundary between two levels of description, the
present-observable level and the predictive level, and the discipline that produced them cuts
both ways. A passing certificate suite never proves emergence, and none of the information-theoretic
emergence criteria in the literature \cite{hoel-2013,rosas-2020} is invoked: GT3's foreclosure says that the
declared lens class provably cannot express the distinction that memory carries, which is the
method working as intended, not a metaphysical discovery about glass, and not a reduction of
either level to the other. Symmetrically, nothing here licenses reading the strictness of the
memory level (GT6) as autonomy, drive, or an arrow of time: no P6\_drive audit was performed, and
a positive closure deficit is neither novelty nor irreversibility on its own. The paper makes no
claim about the existence or non-existence of a thermodynamic ideal-glass transition, takes no
side in the field's thermodynamic-versus-dynamical debate \cite{biroli-garrahan-2013}, and makes
no continuum-limit claims; every quantifier in every claim ranges over a declared finite class, and
the empirical legs are bridge claims only.

\subsection{Full nonclaims ledger}

Per the project's own publication gate, the complete nonclaims ledger, every nonclaim of every
claim row, unabridged, is printed here rather than summarized.

\begin{longtable}{lp{9cm}}
\caption{Full nonclaims ledger (printed in full per the Phase 5 gate)}\label{tab:nonclaims-full} \\
\toprule
claim & nonclaim \\
\midrule
\endfirsthead
\toprule
claim & nonclaim \\
\midrule
\endhead
\bottomrule
\endfoot
GT1 & GT1 certificate is scoped to East N=8 and L\_energy only; FA and L\_panel are undone extensions. \\
GT1 & The tau ladder is the declared finite ladder [1, 2, 4, 8], not the full combinatorial grid. \\
GT1 & Constrained East equilibrium under L\_energy is a nonzero stationary CD baseline; only the unconstrained lumpable control is exact-zero. \\
GT1 & [B] benchmark-parity control is filed as skipped in this artifact because the paper-specific example chains were not ported in this remediation packet. \\
GT2 & filed as P4 staged dependence, not a directionality or irreversibility claim \\
GT2 & Phase-0 Kovacs search output is a pre-certificate diagnostic, not a GT witness certificate. \\
GT2 & The selected result freezes a later config but does not by itself prove a theorem-grade claim. \\
GT3 & finite N=8 East foreclosure over Def(l\_energy), not a continuum or all-lens claim \\
GT4 & GT4 triage is over the declared finite benchmark family, not a shell-general theorem. \\
GT4 & The artifact promotes the Phase-1 benchmark-suite data; it does not change the underlying classification logic. \\
GT4 & GB1 protocol-trap caveat remains a theorem-wrapper limitation; this table files the deterministic triage computation. \\
GT4 & Phase 1 benchmark parity is a control suite, not a standalone GT4 triage certificate. \\
GT4 & Dedicated GT4 regime-table promotion, if present, is filed in a separate artifact. \\
GT5 & buyback saturation or nonsaturation is not a GB5 repair-impossibility certificate; GB5 requires a separate exhaustive repair-class sweep \\
GT5 & the shipped Kovacs buyback curve is a minimal two-history witness with saturation\_budget=1; it demonstrates buyback mechanics, not general single-scalar or TNM insufficiency \\
GT6 & GT6 is scoped to the audited glass shell, not to all East or glassy regimes. \\
GT6 & No shell-general strict-bridge theorem is claimed. \\
GT6 & No broader external-class theorem is claimed beyond the frozen finite class. \\
GT6 & Strictness does not imply autonomous macro closure. \\
GT6 & Strictness does not imply drive; P6\_drive requires the separate GB1 audit discipline. \\
GT6 & The pressure gap is not a KL closure deficit unless an explicit bridge is supplied. \\
GT6 & The pressure gap is not the strictness witness; GT2/GT3 provide the strictness witnesses. \\
GT6 & The gap functional is a conditional disintegration on strata, not stratumwise root separation. \\
GT6 & This is a model realization, not a universal strict-extension theory. \\
GT6 & T1 is a packaged completion layer, not repeated application of one fixed idempotent completion. \\
GT6 & Single bridge only; no stacked or glued promotions are claimed. \\
GT6 & strict-extension inputs are shell-local and do not claim continuum or all-protocol closure \\
GT6 & H5 and non-factorization are cited from prior artifacts, not recomputed here \\
GT6 & Knockout rows are shell-local diagnostics, not shell-general strict-extension proofs. \\
GT6 & P4-P6 rows are capability-loss degradations, not numeric re-simulation degradations. \\
GT6 & Pressure closure is a float64 Fekete diagnostic, not an exact-rational certificate. \\
GT6 & N=10 dense pressure rows may be skipped by declared threshold and are reported as such. \\
GT6 & Spectral gaps are float64 diagnostics and not theorem-graded exact-rational certificates. \\
GT6 & Mixing-time bounds are used for shell-window calibration only. \\
GT7 & run-only readouts, not a fitted scaling law \\
GT7 & East reference only; no transfer claim without a scored FA/trap target run and a pre-registered prediction. \\
GT7 & FA target run only; scoring is deferred to the frozen transfer scorer. \\
GT7 & memdim\_profile and deficit\_decay transfer are out of scope for this target artifact. \\
GT7 & Trap target run only; scoring is deferred to the frozen transfer scorer. \\
GT7 & Trap uses a declared depth observable, not East/FA spin-count energy. \\
GT7 & transfer scoring is mechanical comparison against pre-registered predictions \\
GT7 & score counts are not a fitted scaling law \\
GB1 & GB1 demo records a pseudo-EPR identity and finite numerical check, not a closed general theorem. \\
GB1 & The open obligation remains filed in the payload and is not discharged by this artifact. \\
GB2 & GB2 demo records finite protocol-CD computations, not a GT claim certificate. \\
GB2 & Windowed and marginalized CD values are float64 diagnostics; exactness is reserved for the lumpability check. \\
GB9 & The real Kovacs numeric split-pair witness is not encoded in Lean. \\
GB9 & No non-vacuous PredictiveWitness, StrictRefinement, or LoopAsymmetry instance is claimed. \\
GB9 & The GlassWitness instance is structural and index-based, not a real numeric package export. \\
GB10 & No empirical-realization claim anywhere; bridges are admissibility filings only. \\
GB10 & Visibility records are per instrument; no cross-instrument merge into an instrument-free statement. \\
FINDING.trap\_epsilon\_invariance & Not a novelty claim: the multiplicative-scalar cancellation is standard trap-model structure once stated; the contribution is stating and testing it for the declared family. \\
FINDING.trap\_epsilon\_invariance & Says nothing about trap dynamics or aging trajectories; only the stationary distribution is epsilon-invariant. \\
FINDING.lean\_collapse & Not a claim that Kovacs memory is unformalizable; only that THIS core's coherence axioms foreclose non-vacuous witness instances. \\
FINDING.lean\_collapse & Does not weaken the exact-rational computational witnesses; it explains why they are filed as computational evidence rather than Lean theorems. \\
NEG.material\_forcing\_absent & Absence of material forcing in this shell is not evidence about other shells or protocols. \\
NEG.gt7\_transfer\_failures & Failure counts are not evidence against the run-only thesis; they are the run-only thesis operating as designed. No post-hoc reinterpretation of failed intervals is permitted. \\
NEG.gb5\_repair\_sweep\_not\_built & 'No admissible repair exists' is never asserted at any grade. \\
NEG.gt2\_loop\_certificate\_not\_built & No loop-asymmetry or P3 route-mismatch claim is made anywhere on GT2's behalf. \\
NEG.gt1\_b\_parity\_control\_open & No [B]-parity is claimed for GT1 until the control is actually run. \\
NEG.lean\_instance\_structural & The Lean track certifies the abstract core's theorems, not glass numerics. \\
\bottomrule
\end{longtable}

